\section{מגבולות עבודה מול GUI וסוכני דפדפן}

ההרצאה מציגה הדגמה של סוכן דפדפן (\eng{Claude} עם תוסף, \eng{Nano Browser})
שמבצע חיפוש ביד2 --- ומדגישה את המחיר הכבד של גישה זו. פרק זה מסכם
חסרונות, אלטרנטיבות וכיוון המעבר לעבודה טקסטואלית.

% GUI agent vs API agent
\begin{figure}[H]
\centering
\begin{tikzpicture}
\begin{axis}[
  ybar,
  bar width=18pt,
  ymin=0, ymax=10,
  ylabel={מדד יחסי (גבוה = פחות יעיל)},
  symbolic x coords={GUI Agent, API/Text Agent},
  xtick=data,
  xticklabel style={font=\small},
  legend style={at={(0.5,-0.22)},anchor=north, legend columns=3, font=\footnotesize},
  width=11cm, height=6cm,
  nodes near coords,
  nodes near coords style={font=\tiny}
]
\addplot[fill=accentamber!70] coordinates {(GUI Agent,8) (API/Text Agent,2)};
\addlegendentry{עלות טוקנים}
\addplot[fill=accentblue!60] coordinates {(GUI Agent,9) (API/Text Agent,3)};
\addlegendentry{מורכבות}
\addplot[fill=accentgreen!60] coordinates {(GUI Agent,7) (API/Text Agent,2)};
\addlegendentry{זמן ביצוע}
\end{axis}
\end{tikzpicture}
\caption{סוכן דפדפן מול סוכן טקסט/\eng{API}: עלות, מורכבות וזמן (אילוסטרציה לפי ההרצאה)}
\label{fig:gui-vs-api}
\end{figure}


\subsection{כיצד פועל סוכן דפדפן?}

הסוכן:
\begin{enumerate}[rightmargin=2em]
  \item מקבל מטרה בשפה טבעית (``חפש דירות ברחוב רוטשילד'').
  \item בונה תוכנית עבודה ומבקש אישור.
  \item פותח דפדפן, מצלם מסך (\eng{screenshot}), מנתח תמונה.
  \item מקליד בשדות, לוחץ כפתורים, חוזר עד להשלמה.
\end{enumerate}

כל צילום מסך וניתוח תמונה צורכים טוקנים רבים. כל טעות \eng{UI} (כיוון
כתיבה עברית, חלון מודאלי) מחייבת ניסיון חוזר.

\subsection{חסרונות שזוהו בשיעור}

סטודנטים והמרצה ציינו:
\begin{itemize}[rightmargin=2em]
  \item \textbf{זמן} --- משימה פשוטה נמשכת דקות ארוכות.
  \item \textbf{עלות} --- תשלום לפי טוקנים; צילומי מסך יקרים במיוחד.
  \item \textbf{אבטחה} --- התחברות לגימייל, סיסמאות, אישורי משתמש.
  \item \textbf{RTL} --- אתרים שכותבים משמאל לימין שוברים הקלדה בעברית.
  \item \textbf{שבירות} --- שינוי בעיצוב האתר שובר את הסוכן.
\end{itemize}

\begin{notebox}
המרצה הזכיר כלי ``הפוך על הפוך'' להחלפת כיוון מקלדת עברית-אנגלית ---
פתרון אנקדוטלי, לא ארכיטקטוני.
\end{notebox}

\subsection{מדוע GUI מפריע לעידן הסוכנים?}

החזון מההרצאה: מערכת הפעלה שמבינה שפה חופשית לא צריכה לחקות לחיצות
עכבר. \eng{GUI} תוכנן לבני אדם --- תפריטים, אייקונים, משוב ויזואלי.
סוכן שעובד דרך \eng{GUI} מבזבז משאבים על סימולציה אנושית במקום לגשת
ישירות ל-\eng{API} או לטקסט.

\subsection{השוואה מעשית}

\agenttable{
\caption{סוכן \eng{GUI} מול סוכן \eng{API}/טקסט}
\label{tab:gui-api}
\begin{tabular}{@{}p{3.5cm}p{4cm}p{4cm}@{}}
\toprule
\textbf{מדד} & \textbf{סוכן דפדפן} & \textbf{סוכן טקסט/\eng{API}} \\
\midrule
קלט לסוכן & תמונות מסך & JSON, טקסט, קוד \\
עלות טיפוסית & גבוהה & נמוכה \\
מהירות & איטית & מהירה \\
אמינות & שבירה בעיצוב & יציבה יותר \\
דוגמה במטלה 04 & עריכת Word ב-\eng{GUI} & \eng{LuaLaTeX} + \eng{CLS} \\
\bottomrule
\end{tabular}
}

\subsection{מתי בכל זאת להשתמש ב-GUI Agent?}

יש מקרים שבהם אין \eng{API}:
\begin{itemize}[rightmargin=2em]
  \item אתרים ישנים ללא ממשק מכונה.
  \item טפסים ממשלתיים חד-פעמיים.
  \item משימות חד-פעמיות שלא שוות פיתוח \eng{API}.
\end{itemize}

אך עבור תהליכים \textbf{רפטביליים} --- דוחות, הצעות מחיר, מסמכי קורס ---
יש להעדיף \eng{LaTeX}, \eng{Markdown} ו-\eng{API}.

\begin{insightbox}
המעבר מ-\eng{GUI} לטקסט הוא לא רק טכני --- הוא ארגוני: חברות שמחשיפות
\eng{API} ו-\eng{MCP} מאפשרות אוטומציה; אלה שנשענות רק על ממשק גרפי
נשארות מחוץ למערכת הסוכנים.
\end{insightbox}

\subsection{סיכום פרק}

סוכן דפדפן הוא כלי עוצמתי למשימות אד-הוק, אך לא אדריכלות ברירת מחדל.
למסמכים אקדמיים, דוחות \eng{CTI} ותהליכי ארגון --- עבודה בטקסט עם תבנית
\eng{CLS} עדיפה, זולה וניתנת לבקרת איכות.
